% CV CDI ML / AI Engineer - Stephane KPOVIESSI
% Two-column layout adapted from the freelance CV.
% Compile target: pdfLaTeX

\documentclass[10pt,a4paper]{article}

\usepackage[a4paper,left=0.95cm,right=0.95cm,top=0.55cm,bottom=0.55cm]{geometry}
\usepackage[utf8]{inputenc}
\usepackage[T1]{fontenc}
\usepackage{lmodern}
\usepackage{xcolor}
\usepackage{enumitem}
\usepackage[hidelinks]{hyperref}
\usepackage{tabularx}
\usepackage{titlesec}
\usepackage{fontawesome5}

\definecolor{accent}{HTML}{1F4E79}
\definecolor{accentdark}{HTML}{163A5A}
\definecolor{text}{HTML}{2F3437}
\definecolor{muted}{HTML}{666D70}
\definecolor{rule}{HTML}{D5E2EC}
\definecolor{sidebg}{HTML}{F2F7FB}

\pagestyle{empty}
\setlength{\parindent}{0pt}
\setlength{\parskip}{3.6pt}
\setlist[itemize]{leftmargin=*, topsep=3pt, itemsep=3pt, parsep=0pt, partopsep=0pt}

\titleformat{\section}
  {\fontsize{11.2pt}{13pt}\selectfont\bfseries\color{accent}}
  {}{0pt}{}
  [\vspace{2pt}{\color{rule}\titlerule}\vspace{3pt}]
\titlespacing*{\section}{0pt}{10pt}{5pt}

\newcommand{\smallsection}[2]{%
  \vspace{9pt}
  {\fontsize{8.9pt}{10.45pt}\selectfont\bfseries\color{accentdark}%
  {\scriptsize\color{accent}#1}\hspace{3.5pt}#2}\par
  \vspace{3.2pt}
}

\newcommand{\contactline}[2]{%
  \noindent\makebox[0.34cm][l]{{\color{accent}\scriptsize#1}}%
  \begin{minipage}[t]{\dimexpr\linewidth-0.36cm\relax}%
  \fontsize{8.15pt}{9.5pt}\selectfont #2%
  \end{minipage}\par
  \vspace{1pt}
}

\newcommand{\sideitem}[2]{%
  \textbf{#1}\par
  #2\par
  \vspace{5.7pt}
}

\newcommand{\certitem}[2]{%
  \textbf{#1}\par
  {\fontsize{7.15pt}{8.1pt}\selectfont\color{muted}#2}\par
  \vspace{4.4pt}
}

\newcommand{\role}[3]{%
  \vspace{3pt}
  \begin{tabularx}{\linewidth}{@{}Xr@{}}
  \textbf{#1} & {\color{muted}#2}
  \end{tabularx}\vspace{0.8pt}\\
  {\color{muted}#3}\par
  \vspace{3pt}
}

\newcommand{\projectline}[2]{%
  \textbf{#1} - #2\par
}
\newcommand{\githublink}[2]{%
  \par{\footnotesize\color{accent}\faIcon{github}\ \href{#1}{#2}}\par
}

\newsavebox{\maincolumnbox}
\newlength{\maincolumnheight}

\begin{document}
\fontsize{10pt}{12.3pt}\selectfont
\sloppy

\begin{minipage}[t]{0.52\textwidth}
{\LARGE\bfseries\color{accentdark} Stéphane KPOVIESSI}\par
\vspace{1pt}
{\large\color{accent}Machine Learning / AI Engineer}\par
\vspace{1pt}
{\color{muted}Disponible en CDI}
\end{minipage}
\hfill
\begin{minipage}[t]{0.42\textwidth}
\raggedleft
\vspace{2pt}
{\normalsize\bfseries\color{accent}Machine Learning \textbullet{} Agents IA \textbullet{} RAG}\par
\vspace{3pt}
{\small\color{muted}
Des modèles aux applications IA intégrées aux usages métier.
}
\end{minipage}

\vspace{6pt}
{\color{rule}\hrule height 1pt}
\vspace{7pt}

\begin{lrbox}{\maincolumnbox}
\begin{minipage}[t]{0.695\textwidth}
\vspace{0pt}
\raggedright

\section{Profil}
Ingénieur diplômé en Big Data \& Data Science, je conçois des solutions ML/IA de la préparation des données à l'intégration applicative. Expérience en vision par ordinateur et modèles génératifs chez Sylvagreg, en ingénierie Data chez Allianz, et en développement d'agents IA et de RAG hybride avec Seshat.

\section{Expériences}
\role{Seshat - Fondateur \& AI Engineer}{Sept. 2026 - Aujourd'hui}{Plateforme open source d'IA agentique et d'automatisation}
\begin{itemize}
  \item Conception de services Python/Go, backend local/cloud et application desktop Electron/React pour des agents et workflows métier.
  \item RAG hybride : ingestion documentaire, chunking, embeddings, recherche PostgreSQL/pgvector et Chroma, restitution sourcée.
  \item Orchestration d'agents avec tool calling, MCP, human-in-the-loop et tâches planifiées ; intégration de plusieurs fournisseurs LLM et contrôle des accès.
\end{itemize}
{\footnotesize\color{accent}\faIcon{globe}\ \href{https://seshat-ai.com}{seshat-ai.com}\enspace\textbullet\enspace\faIcon{github}\ \href{https://github.com/KPO-Tech/seshat}{github.com/KPO-Tech/seshat}}\par

\vspace{8pt}
\role{Allianz France - Ingénieur Data / Data Scientist}{Nov. 2025 - Avr. 2026}{Direction P4D, Advanced Data \& Climate - Paris}
\begin{itemize}
  \item Migration d'un datamart Assurance de SAS vers Python/PySpark, avec validation sur plus de 15 000 lignes de code ; architecture Médaillon sur Azure Databricks intégrant plus de 150 règles métier.
  \item Pipeline de Data Quality géospatiale et tests automatisés SAS/Python pour valider les KPI et sécuriser la non-régression.
\end{itemize}

\vspace{8pt}
\role{Groupe Sylvagreg - Machine Learning Engineer}{Mai 2024 - Août 2024}{Lille, France}
\begin{itemize}
  \item R\&D d'une solution de reconstruction 3D à partir d'images 2D, combinant vision par ordinateur, photogrammétrie et modèles génératifs sous PyTorch.
  \item Automatisation des pipelines de prétraitement, augmentation de données et inférence pour fiabiliser les expérimentations et la génération 3D.
  \item Conteneurisation Docker, déploiement AWS et interface de visualisation pour inspecter les reconstructions et faciliter les tests utilisateurs.
\end{itemize}
{\fontsize{7.5pt}{9pt}\selectfont\color{accent}\faIcon{github}\ \href{https://github.com/EngineerProjects/Sylva3dGUI}{github.com/EngineerProjects/Sylva3dGUI}\enspace\textbullet\enspace\href{https://github.com/EngineerProjects/sylva3D}{github.com/EngineerProjects/sylva3D}}\par

\vspace{8pt}
\role{LYNE ET FRERES SARL - Data Analyst}{2023 - 2025}{Missions d'appui stratégique - Bénin / À distance}
\begin{itemize}
  \item Analyse de données métiers avec Python/Pandas et restitution via rapports et interfaces Streamlit / web.
  \item Déploiement d'automatisations et d'outils IA internes open source (OpenWebUI, modèles ouverts), avec accompagnement des employés.
\end{itemize}

\section{Projets sélectionnés}
\projectline{SOLAR IRRADIANCE FORECASTING}{Scikit-learn, Pandas, FastAPI}
Préparation de données météo, feature engineering temporel et évaluation de modèles de prévision ; API REST FastAPI pour exposer les résultats.
\githublink{https://github.com/EngineerProjects/solar-forecasting}{github.com/EngineerProjects/solar-forecasting}

\vspace{9pt}
\projectline{ROBOT AUTONOME DE COLLECTE DE DÉCHETS}{AceObject - YOLOv8, OpenCV, Dijkstra}
Pipeline Python/YOLOv8/OpenCV pour détecter, classifier et suivre les déchets en les distinguant des obstacles. Navigation par Dijkstra pour planifier la collecte ; exploration de l'apprentissage par renforcement pour la décision et l'évitement d'obstacles.

\end{minipage}
\end{lrbox}
\setlength{\maincolumnheight}{\dimexpr\ht\maincolumnbox+\dp\maincolumnbox\relax}

\begin{minipage}[t]{0.285\textwidth}
\vspace{0pt}

\hspace*{-0.18cm}%
\colorbox{sidebg}{%
  \hspace*{0.18cm}%
  \begin{minipage}[t][\maincolumnheight][t]{0.91\linewidth}
  \vspace{8pt}
  \fontsize{8.45pt}{10.35pt}\selectfont
  \raggedright

  \smallsection{\faIcon{address-card}}{Contact}
  \contactline{\faIcon{map-marker-alt}}{Cergy / Île-de-France}
  \contactline{\faIcon{phone}}{+33 7 48 59 47 13}
  \contactline{\faIcon{envelope}}{\href{mailto:oastephanekpoviessi@gmail.com}{oastephanekpoviessi@gmail.com}}
  \contactline{\faIcon{link}}{\href{https://linkedin.com/in/stephanekpoviessi}{linkedin.com/in/stephanekpoviessi}}
  \contactline{\faIcon{github}}{\href{https://github.com/KPO-Tech}{github.com/KPO-Tech}}
  \contactline{\faIcon{globe}}{\href{https://seshat-ai.com}{seshat-ai.com}}

  \smallsection{\faIcon{layer-group}}{Compétences}
  \sideitem{ML \& Computer Vision}{PyTorch, Scikit-learn, prétraitement, augmentation, évaluation, modèles génératifs.}
  \sideitem{Agents \& RAG}{LLM, LangGraph, MCP, tool calling, embeddings, recherche hybride.}
  \sideitem{Data \& stockage}{Python, PySpark, Pandas, SQL, PostgreSQL/pgvector, Chroma.}
  \sideitem{Engineering}{FastAPI, Go, Azure Databricks, AWS, Docker, Git/CI-CD.}

  \smallsection{\faIcon{graduation-cap}}{Formation}
  \sideitem{JUNIA ISEN Lille}{Diplôme d'Ingénieur - Big Data \& Data Science, 2026.}
  \sideitem{JUNIA ISEN Rabat}{Licence - Électronique \& Systèmes Numériques, 2023.}

  \smallsection{\faIcon{certificate}}{Certifications}
  \certitem{\href{https://learn.microsoft.com/fr-fr/training/career-paths/ai-engineer}{Microsoft AI Engineer Career Path}}{Microsoft Learn \textbullet{} \textbf{En cours}}
  \certitem{\href{https://www.coursera.org/account/accomplishments/verify/LU2ZQ972BWNF}{Create ML Models in Microsoft Azure}}{Microsoft / Coursera \textbullet{} 2023}
  \certitem{\href{https://www.coursera.org/account/accomplishments/verify/MWJS4Q3S2C3L}{Supervised ML: Regression \& Classification}}{DeepLearning.AI / Stanford Online\\2023}
  \certitem{\href{https://www.freecodecamp.org/certification/amiche/data-analysis-with-python-v7}{Data Analysis with Python}}{freeCodeCamp \textbullet{} 2023}

  \smallsection{\faIcon{language}}{Langues}
  \textbf{Français} - langue maternelle\par
  \textbf{Anglais} - B2, Linguaskill Business\par

  \smallsection{\faIcon{users}}{{\fontsize{8.4pt}{10pt}\selectfont Engagements \& associations}}
  \textbf{JUNIA ALUMNI}\par
  {\color{muted}Commission Réseau}\par
  \vspace{3pt}
  \textbf{Vice-président}\par
  {\color{muted}Sept. 2024 - Août 2025}\par
  Animation du réseau alumni et événements professionnels.\par
  \vspace{4pt}
  \textbf{Responsable communication}\par
  {\color{muted}Sept. 2023 - Août 2024}\par
  Annonces et valorisation des initiatives.\par
  \vspace{3pt}
  \end{minipage}
}

\end{minipage}
\hfill
\usebox{\maincolumnbox}

\end{document}
